\section*{Chapter \ref{ch:g6:latitude-longitude} --- Latitude, Longitude and Time Zones}

\begin{solution}{exo:g6:latitude-longitude:1}
``\omterm{def:g6:latitude-longitude:latitude}{Latitude} is measured from the \omterm{def:g3:continents-and-oceans:equator}{equator} and goes from $0^\circ$ to
$90^\circ$ (north or south); \omterm{def:g6:latitude-longitude:longitude}{longitude} is measured from the \omterm{def:g6:latitude-longitude:longitude}{prime meridian}
and goes from $0^\circ$ to $180^\circ$ (east or west).''
\end{solution}

\begin{solution}{exo:g6:latitude-longitude:2}
The \omterm{def:g3:continents-and-oceans:poles}{North Pole}: $90^\circ$\,N. The \omterm{def:g3:continents-and-oceans:poles}{South Pole}: $90^\circ$\,S. A town on the
\omterm{def:g3:continents-and-oceans:equator}{equator}: $0^\circ$.
\end{solution}

\begin{solution}{exo:g6:latitude-longitude:3}
Sydney: \omterm{def:g3:continents-and-oceans:equator}{southern hemisphere}, east. New York: \omterm{def:g3:continents-and-oceans:equator}{northern hemisphere}, west.
Cairo: \omterm{def:g3:continents-and-oceans:equator}{northern hemisphere}, east. Rio de Janeiro: \omterm{def:g3:continents-and-oceans:equator}{southern hemisphere},
west.
\end{solution}

\begin{solution}{exo:g6:latitude-longitude:4}
(a) True. (b) False: the \omterm{def:g6:latitude-longitude:latitude}{parallels} get shorter towards the \omterm{def:g3:continents-and-oceans:poles}{poles}; the
longest is the \omterm{def:g3:continents-and-oceans:equator}{equator}. (c) True. (d) False: it passes through Greenwich,
in London.
\end{solution}

\begin{solution}{exo:g6:latitude-longitude:5}
Cairo. New York. Tokyo is at about $36^\circ$\,N, $140^\circ$\,E (a little
above the \omterm{def:g6:latitude-longitude:latitude}{parallel} $30^\circ$\,N, a little beyond the \omterm{def:g6:latitude-longitude:longitude}{meridian}
$120^\circ$\,E, short of $150^\circ$\,E).
\end{solution}

\begin{solution}{exo:g6:latitude-longitude:6}
$45 \div 15 = 3$: their noons are 3~hours apart. It is noon first in
town~B, which is further east.
\end{solution}

\begin{solution}{exo:g6:latitude-longitude:7}
$73 \div 15 \approx 4.9$, rounded to 5: the zone is $+5$. When it is noon
at Greenwich, it is $12 + 5 = 17$, that is 5~p.m., in Mumbai.
\end{solution}

\begin{solution}{exo:g6:latitude-longitude:8}
$30 \times 111 = 3\,330$~km. The \omterm{def:g3:continents-and-oceans:poles}{North Pole} is at $90^\circ$:
$90 \times 111 = 9\,990$~km, about 10\,000~km.
\end{solution}

\begin{solution}{exo:g6:latitude-longitude:9}
$9 - (-5) = 14$~hours apart. $8 + 14 = 22$: it is 10~p.m. on Monday in
Tokyo.
\end{solution}

\begin{solution}{exo:g6:latitude-longitude:10}
She is right. The Earth turns from west to east, so the towns in the east
are carried towards the Sun first: they see it rise, reach its highest
point and set before the towns further west.
\end{solution}

\begin{solution}{exo:g6:latitude-longitude:11}
(a) $42 \times 3 = 126$~seconds, a little more than 2~minutes. (b)
$126 \div 240 \approx 0.5$: about half a degree of \omterm{def:g6:latitude-longitude:longitude}{longitude}. (c)
$0.5 \times 111 \approx 55$~km (with $0.525$: about 58~km). A mistake of
about fifty kilometres: enough to miss an island, or to run onto a \omterm{def:g1:landscapes-around-us:water}{coast}
thought to be far away.
\end{solution}

\begin{solution}{exo:g6:latitude-longitude:12}
(a) $18 - 9 = 9$: 9~a.m. on Monday at Greenwich. (b) $9 + 10 = 19$:
7~p.m. on Monday at Greenwich. (c) $19 - 8 = 11$: 11~a.m. on Monday in Los
Angeles. The plane left at 6~p.m. on Monday (Tokyo time) and lands at
11~a.m. on the same Monday (Los Angeles time): the clocks at arrival show an
earlier time than the clocks at departure, because the plane crossed the
date line towards the east and went back one day.
\end{solution}

\begin{solution}{pb:g6:latitude-longitude:1}
\textbf{1.} Ana (Rio de Janeiro, $23^\circ$\,S) and Mia (Auckland,
$37^\circ$\,S).

\textbf{2.} Ana and Sam ($43^\circ$\,W and $118^\circ$\,W). Lena is on the
\omterm{def:g6:latitude-longitude:longitude}{prime meridian}.

\textbf{3.} Lena: London is at $52^\circ$\,N, farther from the \omterm{def:g3:continents-and-oceans:equator}{equator}
than all the others.

\textbf{4.} Mia: Auckland is at $175^\circ$\,E, only $180 - 175 = 5$
degrees from the $180^\circ$ \omterm{def:g6:latitude-longitude:longitude}{meridian}.

\textbf{5.} $360^\circ \div 24 = 15^\circ$.

\textbf{6.} London: $0$. Rio de Janeiro: $43 \div 15 \approx 2.9$, west,
so $-3$. Tokyo: $140 \div 15 \approx 9.3$, east, so $+9$. Auckland:
$175 \div 15 \approx 11.7$, east, so $+12$. Los Angeles:
$118 \div 15 \approx 7.9$, west, so $-8$.

\textbf{7.} Rio de Janeiro: $20 - 3 = 17$, 5~p.m. on Friday. Tokyo:
$20 + 9 = 29$; $29 - 24 = 5$: 5~a.m. on Saturday.

\textbf{8.} Tokyo is far to the east: its clocks are 9~hours ahead of
London's. At 8~p.m. in London, Tokyo has already passed midnight, so its
Saturday has begun.

\textbf{9.} $20 - 8 = 12$: noon on Friday.

\textbf{10.} $12 - (-8) = 20$~hours ahead.

\textbf{11.} Crossing the date line towards the east, the ship goes back to
the day before.

\textbf{12.} $20 + 12 = 32$; $32 - 24 = 8$: when the call starts, it is
\textbf{8~a.m. on Saturday} in Auckland. Mia joins the call at breakfast
time, while Sam has his lunch on Friday.
\end{solution}
